
%%%%%%%%%%%%%%%%%%%%%%%%%%%%%%%%%%%%%%%%%%%%%%%%%%%%%%%%%%%%%%%%%%%
% Physik 4 - Blatt 10 Abgabe
%%%%%%%%%%%%%%%%%%%%%%%%%%%%%%%%%%%%%%%%%%%%%%%%%%%%%%%%%%%%%%%%%%%

\documentclass[11pt,a4paper]{scrartcl}

% Layout
\usepackage[margin=2.5cm]{geometry}
\usepackage[onehalfspacing]{setspace}

% General packages
\usepackage{graphicx}
\usepackage{enumitem}
\usepackage{tcolorbox}
\tcbuselibrary{breakable}
\usepackage[breaklinks=true,colorlinks=true,linkcolor=blue,urlcolor=blue,citecolor=blue]{hyperref}
\usepackage{amsmath,amsthm,amssymb,mathtools}
\usepackage{icomma}
\usepackage[T1]{fontenc}
\usepackage[utf8]{inputenc}
\usepackage[ngerman]{babel}
\usepackage{tikz}

% Units / tables / floats / references
\usepackage[locale = DE,space-before-unit=true,per-mode = symbol]{siunitx}
\usepackage{booktabs,multirow}
\usepackage{placeins}
\usepackage[natbib,abbreviate=true,doi=false,style=numeric-comp,giveninits=true,sorting=none]{biblatex}
\usepackage{csquotes}
\usepackage{pgfplots}
\pgfplotsset{compat=1.18}

% Graphics paths
\graphicspath{{Bilder/} {images/} {figures/} {chapters/}}

% Optional math shortcuts
\newcommand{\R}{\mathbb{R}}
\newcommand{\N}{\mathbb{N}}
\newcommand{\Q}{\mathbb{Q}}
\newcommand{\set}[1]{\left\{ #1 \right\}}
\newcommand{\abs}[1]{\left| #1 \right|}
\newcommand{\norm}[1]{\left\lVert #1 \right\rVert}
\newcommand{\ol}[1]{\overline{#1}}
\newcommand{\half}{\frac{1}{2}}
\newcommand{\dd}{\,\mathrm{d}}

% Optional units
\DeclareSIUnit{\year}{a}
\DeclareSIUnit{\day}{d}

% Box styles
\newtcolorbox{calculationbox}[1][]{
    title=Calculation,
    breakable,
    #1
}

\newtcolorbox{solutionbox}[1][]{
    title=Solution,
    breakable,
    #1
}

\newtcolorbox{answerbox}[1][]{
    title=Answer,
    breakable,
    #1
}

\newtcolorbox{notebox}[1][]{
    title=Notes,
    breakable,
    #1
}

\newtcolorbox{sidecalcbox}[1][]{
    title=Side Calculation,
    breakable,
    #1
}

%%%%%%%%%%%%%%%%%%%%%%%%%%%%%%%%%%%%%%%%%%%%%%%%%%%%%%%%%%%%%%%%%%%
\begin{document}

    \title{Physik 4 - Blatt 10}
    \author{Tobias Laser}
    \date{\today}
    \maketitle
    \vfill
    \thispagestyle{empty}
    \newpage

    \pagenumbering{arabic}

    \paragraph{Aufgabe 1: Elektromagnetische Teilchenschauer}

    Teilchenschauer sind u. a. für die Beobachtung von hochenergetischen Teilchen der
    kosmischen Strahlung relevant. In diesem Fall wird deren Wechselwirkung mit der
    Erdatmosphäre untersucht. Im Folgenden werden wir der Einfachheit halber nur
    elektromagnetische Schauer betrachten.

    In einem elektromagnetischen Schauer sind die dominierenden Prozesse Bremsstrahlung
    und Paarbildung. Wir betrachten hier ein vereinfachtes Modell für einen solchen
    Schauer nach dem deutschen Physiker Walter Heinrich Heitler. Dabei nimmt man an,
    daß ein hochenergetisches Elektron mit Energie $E_0$, das auf die Erdatmosphäre
    trifft, nach einer Strahlungslänge $X_0$ ein Photon der Energie $E_0/2$ abstrahlt.
    Nach einer weiteren Strecke $X_0$ wird das Photon in ein Elektron-Positron-Paar
    umgewandelt, wobei die Energie ebenfalls gleichmäßig aufgeteilt wird. Ebenso strahlt
    das primäre Elektron unter Abgabe seiner halben Energie ein weiteres Photon ab.
    Auf diese Weise entwickelt sich der Schauer so lange weiter, bis die Teilchen eine
    kritische Energie $E_c$ erreichen, ab der keine neuen Teilchen mehr produziert werden
    können. Nach diesem Schauermaximum verlieren die Teilchen durch verschiedene Prozesse
    ihre Restenergie und können bald nicht mehr beobachtet werden.

    \begin{enumerate}[label=(\alph*)]
        \item
            Skizzieren Sie zuerst die ersten drei Stufen des Schauers und nutzen Sie
            unterschiedliche Symbole für Leptonen und Photonen.

            \begin{calculationbox}
                Wir unterscheiden Leptonen, also Elektronen und Positronen, von Photonen.
                Eine einfache symbolische Wahl ist
                \[
                    e^-,e^+ \text{ als durchgezogene Linien},
                    \qquad
                    \gamma \text{ als gewellte Linien}.
                \]
                Die ersten drei Stufen des Heitler-Modells lassen sich dann so darstellen:

                \begin{center}
                    \begin{tikzpicture}[scale=1.0,>=stealth]
                        \foreach \x/\lab in {0/{0},2/{$X_0$},4/{$2X_0$},6/{$3X_0$}}{
                            \draw[dashed,gray] (\x,-2.0) -- (\x,2.0);
                            \node[above] at (\x,2.0) {\lab};
                        }

                        \draw[->,thick] (0,0) -- (2,0.9) node[midway,above] {$e^-$};
                        \draw[decorate,decoration={snake,amplitude=1.5pt,segment length=6pt},->,thick]
                            (2,0.9) -- (4,1.6) node[midway,above] {$\gamma$};
                        \draw[->,thick] (2,0.9) -- (4,0.55) node[midway,below] {$e^-$};

                        \draw[->,thick] (4,1.6) -- (6,1.85) node[midway,above] {$e^-$};
                        \draw[->,thick] (4,1.6) -- (6,1.25) node[midway,below] {$e^+$};
                        \draw[->,thick] (4,0.55) -- (6,0.20) node[midway,below] {$e^-$};
                        \draw[decorate,decoration={snake,amplitude=1.5pt,segment length=6pt},->,thick]
                            (4,0.55) -- (6,0.85) node[midway,above] {$\gamma$};

                        \node[left] at (0,0) {$e^-(E_0)$};
                        \node[below] at (0,-2.0) {Start};
                    \end{tikzpicture}
                \end{center}

                \begin{calculationbox}
                    Die Skizze ist nicht als exakte Feynman-Darstellung gemeint, sondern als
                    Baumdiagramm des vereinfachten Heitler-Modells: Nach jeder Strahlungslänge
                    verdoppelt sich idealisiert die Teilchenzahl und die Energie pro Teilchen
                    halbiert sich.
                \end{calculationbox}
            \end{calculationbox}

        \item
            Finden Sie die Strahlungslänge $X_{\max}$, bei der das Schauermaximum erreicht
            wird. Geben Sie das Ergebnis in Einheiten von $X_0$ an. Wie groß ist die Zahl
            $N_{\max}$ der Teilchen im Schauermaximum?

            \begin{calculationbox}
                Im Heitler-Modell verdoppelt sich nach jeder Stufe die Teilchenzahl. Nach $n$
                Stufen gilt daher
                \[
                    N(n)=2^n.
                \]
                Gleichzeitig wird die Energie gleichmäßig auf alle Teilchen verteilt:
                \[
                    E(n)=\frac{E_0}{2^n}.
                \]

                \begin{calculationbox}
                    Das Schauermaximum ist erreicht, wenn die Energie pro Teilchen gerade auf
                    die kritische Energie abgesunken ist:
                    \[
                        E(n_{\max})=E_c.
                    \]
                    Also
                    \begin{align*}
                        \frac{E_0}{2^{n_{\max}}}&=E_c,\
                        2^{n_{\max}}&=\frac{E_0}{E_c},\
                        n_{\max}&=\frac{\ln(E_0/E_c)}{\ln 2}.
                    \end{align*}
                \end{calculationbox}

                \begin{calculationbox}
                    Da jede Stufe einer Strahlungslänge $X_0$ entspricht, folgt
                    \[
                        X_{\max}=n_{\max}X_0
                        =X_0\frac{\ln(E_0/E_c)}{\ln 2}.
                    \]
                    In Einheiten von $X_0$ ist also
                    \[
                        \frac{X_{\max}}{X_0}=\frac{\ln(E_0/E_c)}{\ln 2}.
                    \]
                    Die Teilchenzahl im Maximum ist
                    \[
                        N_{\max}=2^{n_{\max}}=\frac{E_0}{E_c}.
                    \]
                \end{calculationbox}
            \end{calculationbox}

            \begin{answerbox}
                Das Schauermaximum liegt bei
                \[
                    \boxed{\frac{X_{\max}}{X_0}=\frac{\ln(E_0/E_c)}{\ln 2}}
                \]
                und die maximale Teilchenzahl ist
                \[
                    \boxed{N_{\max}=\frac{E_0}{E_c}.}
                \]
            \end{answerbox}

        \item
            Für elektromagnetische Schauer in Luft beträgt $E_c$ circa \SI{84}{MeV}.
            Welche Schauerlängen in Einheiten von $X_0$ ergeben sich damit für primäre
            Photonen mit einer Energie von $10^{11}\,\si{eV}$ bzw. $10^{16}\,\si{eV}$?
            Gehen Sie dazu davon aus, daß das Schauermaximum etwa bei $2/3$ der
            Schauerlänge liegt. Welche physikalischen Prozesse determinieren die kritische
            Energie?

            \begin{calculationbox}
                Zunächst rechnen wir die kritische Energie in eV um:
                \[
                    E_c=\SI{84}{MeV}=8.4\cdot 10^7\,\si{eV}.
                \]
                Für das Schauermaximum verwenden wir aus Teil (b)
                \[
                    \frac{X_{\max}}{X_0}=\frac{\ln(E_0/E_c)}{\ln 2}.
                \]
                Wenn das Maximum bei etwa $2/3$ der gesamten Schauerlänge $L$ liegt, gilt
                \[
                    X_{\max}\approx \frac{2}{3}L
                    \qquad\Longrightarrow\qquad
                    L\approx \frac{3}{2}X_{\max}.
                \]

                \begin{calculationbox}
                    Für $E_0=10^{11}\,\si{eV}$ ergibt sich
                    \begin{align*}
                        \frac{X_{\max}}{X_0}
                        &=\frac{\ln(10^{11}/(8.4\cdot10^7))}{\ln2}\\
                        &=10.22.
                    \end{align*}
                    Damit ist die gesamte Schauerlänge näherungsweise
                    \[
                        \frac{L}{X_0}\approx\frac{3}{2}\cdot 10.22=15.33.
                    \]
                \end{calculationbox}

                \begin{calculationbox}
                    Für $E_0=10^{16}\,\si{eV}$ ergibt sich analog
                    \begin{align*}
                        \frac{X_{\max}}{X_0}
                        &=\frac{\ln(10^{16}/(8.4\cdot10^7))}{\ln2}\\
                        &=26.83.
                    \end{align*}
                    Damit ist
                    \[
                        \frac{L}{X_0}\approx\frac{3}{2}\cdot 26.83=40.24.
                    \]
                \end{calculationbox}

                \begin{calculationbox}
                    Die kritische Energie wird physikalisch dadurch bestimmt, daß bei dieser
                    Energie die dominierenden Verlustmechanismen wechseln. Für Elektronen und
                    Positronen ist $E_c$ näherungsweise die Energie, bei der Strahlungsverluste
                    durch Bremsstrahlung und kontinuierliche Energieverluste, vor allem durch
                    Ionisation und Anregung des Mediums, vergleichbar groß werden. Unterhalb
                    von $E_c$ ist die weitere Teilchenproduktion ineffizient.
                \end{calculationbox}
            \end{calculationbox}

            \begin{answerbox}
                Für $10^{11}\,\si{eV}$ erhält man
                \[
                    \boxed{\frac{X_{\max}}{X_0}=10.22,\qquad \frac{L}{X_0}=15.33.}
                \]
                Für $10^{16}\,\si{eV}$ erhält man
                \[
                    \boxed{\frac{X_{\max}}{X_0}=26.83,\qquad \frac{L}{X_0}=40.24.}
                \]
            \end{answerbox}

        \item
            Die atmosphärische Tiefe $X$ in Einheiten von $\si{g/cm^2}$ hängt exponentiell
            von der Höhe $h$ über dem Meeresspiegel ab:
            \[
                X=X_Ee^{-h/H_0},
            \]
            wobei $H_0=\SI{8}{km}$ und $X_E=\SI{1030}{g/cm^2}$ die totale atmosphärische
            Tiefe ist. In welcher Höhe befindet sich also das Schauermaximum für die Photonen
            aus (c), wenn die Strahlungslänge $X_0=\SI{37.8}{g/cm^2}$ beträgt?

            \begin{calculationbox}
                Aus Teil (c) kennen wir $X_{\max}/X_0$. Daraus erhalten wir die atmosphärische
                Tiefe des Maximums:
                \[
                    X_{\max}^{(\si{g/cm^2})}=\frac{X_{\max}}{X_0}\cdot X_0.
                \]
                Anschließend lösen wir
                \[
                    X=X_Ee^{-h/H_0}
                \]
                nach $h$ auf:
                \[
                    h=-H_0\ln\left(\frac{X}{X_E}\right).
                \]

                \begin{calculationbox}
                    Für $E_0=10^{11}\,\si{eV}$ ist
                    \begin{align*}
                        X_{\max}&=10.22\cdot \SI{37.8}{g/cm^2}\\
                        &=\SI{386.21}{g/cm^2}.
                    \end{align*}
                    Daher
                    \begin{align*}
                        h&=-\SI{8}{km}\cdot\ln\left(\frac{386.21}{1030}\right)\\
                        &=\SI{7.85}{km}.
                    \end{align*}
                \end{calculationbox}

                \begin{calculationbox}
                    Für $E_0=10^{16}\,\si{eV}$ ist
                    \begin{align*}
                        X_{\max}&=26.83\cdot \SI{37.8}{g/cm^2}\\
                        &=\SI{1014.06}{g/cm^2}.
                    \end{align*}
                    Daher
                    \begin{align*}
                        h&=-\SI{8}{km}\cdot\ln\left(\frac{1014.06}{1030}\right)\\
                        &=\SI{0.12}{km}.
                    \end{align*}
                \end{calculationbox}
            \end{calculationbox}

            \begin{answerbox}
                Für ein primäres Photon mit $10^{11}\,\si{eV}$ liegt das Schauermaximum bei
                \[
                    \boxed{h=\SI{7.85}{km}.}
                \]
                Für ein primäres Photon mit $10^{16}\,\si{eV}$ liegt es bei
                \[
                    \boxed{h=\SI{0.12}{km}.}
                \]
            \end{answerbox}

        \item
            Bestimmen Sie schließlich die Energie, die ein Elektron der kosmischen Strahlung
            haben muß, damit das Schauermaximum gerade am Technik-Campus der Universität
            Innsbruck auftritt.

            \begin{calculationbox}
                Für die Höhe des Technik-Campus verwenden wir näherungsweise die Höhe von
                Innsbruck über dem Meeresspiegel,
                \[
                    h\approx\SI{0.574}{km}.
                \]
                Mit der barometrischen Formel ist die atmosphärische Tiefe an dieser Höhe
                \[
                    X(h)=X_Ee^{-h/H_0}.
                \]

                \begin{calculationbox}
                    Einsetzen liefert
                    \begin{align*}
                        X(h)&=\SI{1030}{g/cm^2}\cdot e^{-0.574/8}\\
                        &=\SI{958.69}{g/cm^2}.
                    \end{align*}
                    In Einheiten der Strahlungslänge ist das
                    \begin{align*}
                        \frac{X_{\max}}{X_0}
                        &=\frac{958.69}{37.8}\\
                        &=25.36.
                    \end{align*}
                \end{calculationbox}

                \begin{calculationbox}
                    Aus
                    \[
                        \frac{X_{\max}}{X_0}=\frac{\ln(E_0/E_c)}{\ln2}
                    \]
                    folgt
                    \[
                        E_0=E_c\,2^{X_{\max}/X_0}.
                    \]
                    Damit ergibt sich
                    \begin{align*}
                        E_0
                        &=8.4\cdot10^7\,\si{eV}\cdot 2^{25.36}\\
                        &=3.62\cdot10^{15}\,\si{eV}.
                    \end{align*}
                \end{calculationbox}
            \end{calculationbox}

            \begin{answerbox}
                Ein Elektron der kosmischen Strahlung müßte näherungsweise die Energie
                \[
                    \boxed{E_0\approx 3.62\cdot 10^{15}\,\si{eV}=\SI{3.62}{PeV}}
                \]
                haben, damit das Schauermaximum gerade in etwa auf Höhe des Technik-Campus
                Innsbruck liegt. Das Ergebnis hängt leicht von der konkret verwendeten
                Campus-Höhe ab.
            \end{answerbox}
    \end{enumerate}

    \newpage
    \paragraph{Aufgabe 2: Elektromagnetische Teilchenschauer}

    Die Ausbreitungsgeschwindigkeit von Licht in einem optisch durchlässigen Medium ist
    um dessen Brechungsindex $n$ verringert, also $c_n=c/n$. Durchquert nun ein geladenes
    Teilchen mit Geschwindigkeit $v=\beta c>c_n$ dieses Medium, so kommt es zur Emission
    von Cherenkov-Licht. Dieses wird relativ zur Flugrichtung unter einem Winkel
    $\vartheta_c$ emittiert, welcher sich folgendermaßen berechnen lässt:
    \[
        \cos\vartheta_c=\frac{1}{\beta n}.
    \]
    Der Cherenkov-Effekt tritt unter anderem in der Atmosphäre auf, wenn ein hochenergetisches
    Teilchen der kosmischen Strahlung oder auch ein Gamma-Quant einen Teilchenschauer auslöst.

    \begin{enumerate}[label=(\alph*)]
        \item
            Berechnen Sie $\vartheta_c$ für ein Positron mit $E=\SI{200}{MeV}$ in einer Höhe
            $h=\SI{10}{km}$ ü. NN. Der Brechungsindex in Luft ist dabei über die
            barometrische Höhenformel als
            \[
                n(h)=1+0.000283\cdot e^{-h/H_0},\qquad H_0=\SI{8}{km}
            \]
            gegeben.

            \begin{calculationbox}
                Für ein Positron ist die Ruheenergie
                \[
                    m_ec^2=\SI{0.511}{MeV}.
                \]
                Bei einer Gesamtenergie von $E=\SI{200}{MeV}$ gilt
                \[
                    \gamma=\frac{E}{m_ec^2}.
                \]
                Daraus folgt
                \[
                    \beta=\sqrt{1-\frac{1}{\gamma^2}}.
                \]

                \begin{calculationbox}
                    Numerisch erhalten wir
                    \begin{align*}
                        \gamma&=\frac{200}{0.511}=391.39,\\
                        \beta&=\sqrt{1-\frac{1}{391.39^2}}=0.9999967.
                    \end{align*}
                    Der Brechungsindex in $h=\SI{10}{km}$ Höhe ist
                    \begin{align*}
                        n(10\,\si{km})
                        &=1+0.000283\cdot e^{-10/8}\\
                        &=1.0000811.
                    \end{align*}
                \end{calculationbox}

                \begin{calculationbox}
                    Damit ist
                    \begin{align*}
                        \cos\vartheta_c
                        &=\frac{1}{\beta n}\\
                        &=\frac{1}{0.9999967\cdot1.0000811}\\
                        &=0.9999222.
                    \end{align*}
                    Also
                    \begin{align*}
                        \vartheta_c
                        &=\arccos(0.9999222)\\
                        &=0.01247\,\si{rad}\\
                        &=\SI{0.714}{\degree}.
                    \end{align*}
                \end{calculationbox}
            \end{calculationbox}

            \begin{answerbox}
                Der Cherenkov-Winkel beträgt
                \[
                    \boxed{\vartheta_c=0.01247\,\si{rad}=\SI{0.714}{\degree}.}
                \]
            \end{answerbox}

        \item
            Angenommen das Positron aus der vorherigen Teilaufgabe kommt direkt aus dem Zenit:
            Wie groß ist der Radius des resultierenden Cherenkov-Kegels auf Meereshöhe?

            \begin{calculationbox}
                Wenn das Positron senkrecht nach unten läuft und das Licht in Höhe
                $h=\SI{10}{km}$ unter dem Winkel $\vartheta_c$ emittiert wird, ergibt sich
                für den Radius des Kegels auf Meereshöhe näherungsweise
                \[
                    r=h\tan\vartheta_c.
                \]
                Dabei verwenden wir den in Teil (a) berechneten Winkel.

                \begin{calculationbox}
                    Mit
                    \[
                        h=\SI{10}{km}=\SI{10000}{m},
                        \qquad
                        \vartheta_c=0.01247\,\si{rad}
                    \]
                    folgt
                    \begin{align*}
                        r&=\SI{10000}{m}\cdot \tan(0.01247)\\
                        &=\SI{124.7}{m}.
                    \end{align*}
                \end{calculationbox}
            \end{calculationbox}

            \begin{answerbox}
                Der Radius des Cherenkov-Kegels auf Meereshöhe ist näherungsweise
                \[
                    \boxed{r\approx\SI{125}{m}.}
                \]
            \end{answerbox}
    \end{enumerate}

    \newpage
    \paragraph{Aufgabe 3: Energieverlust durch Ionisation}

    Der Energieverlust von geladenen Teilchen in Materie wird durch die Bethe-Bloch-Formel
    beschrieben:
    \[
        -\frac{\dd E}{\dd x}
        =\frac{1}{(4\pi\epsilon_0)^2}\frac{4\pi e^4z^2}{m_ec^2}
        N_A\frac{Z}{A}\rho\frac{1}{\beta^2}
        \left(
            \ln\left(\frac{2m_ec^2\beta^2\gamma^2}{I}\right)-\beta^2
        \right).
    \]
    Dabei ist $\beta=v/c$, $z$ die Kernladungszahl des Teilchens, $Z$ die des
    Targetmaterials, $\rho$ dessen Dichte, $I$ das mittlere Anregungspotential der Atome
    und $A$ die molare Masse in $\si{g/mol}$.

    \begin{enumerate}[label=(\alph*)]
        \item
            Berechnen und vergleichen Sie den Energieverlust von Protonen, Deuteronen und
            $\alpha$-Teilchen mit kinetischer Energie $\SI{1}{GeV}$ beim Durchdringen eines
            Graphitabsorbers der Dicke $\SI{2}{cm}$. Geben Sie alle für die Berechnung
            notwendigen Eingangsgrößen explizit an.
            \[
                m_p=\SI{938.3}{MeV/c^2},\quad
                m_D=\SI{1875.6}{MeV/c^2},\quad
                m_{\alpha}=\SI{3727.4}{MeV/c^2},
            \]
            \[
                A_C=\SI{12.01}{g/mol},\quad
                Z_C=6,\quad
                \rho_{\mathrm{gr}}=\SI{2.2}{g/cm^3},\quad
                I\approx 16\cdot Z^{0.9}\,\si{eV}.
            \]

            \begin{calculationbox}
                Für die praktische Rechnung verwenden wir die übliche Konstante
                \[
                    K=\frac{1}{(4\pi\epsilon_0)^2}\frac{4\pi e^4}{m_ec^2}N_A
                    \approx \SI{0.307075}{MeV.cm^2/mol}.
                \]
                Dann lautet die verwendete Form
                \[
                    -\frac{\dd E}{\dd x}
                    =Kz^2\frac{Z}{A}\rho\frac{1}{\beta^2}
                    \left(
                        \ln\left(\frac{2m_ec^2\beta^2\gamma^2}{I}\right)-\beta^2
                    \right).
                \]

                \begin{calculationbox}
                    Für Graphit setzen wir ein:
                    \[
                        Z=6,
                        \qquad
                        A=\SI{12.01}{g/mol},
                        \qquad
                        \rho=\SI{2.2}{g/cm^3}.
                    \]
                    Das mittlere Anregungspotential ist
                    \begin{align*}
                        I&=16\cdot 6^{0.9}\,\si{eV}\\
                        &=\SI{80.25}{eV}\\
                        &=8.025\cdot10^{-5}\,\si{MeV}.
                    \end{align*}
                    Außerdem ist
                    \[
                        m_ec^2=\SI{0.511}{MeV},
                        \qquad
                        d=\SI{2}{cm}.
                    \]
                \end{calculationbox}

                \begin{calculationbox}
                    Für jedes Teilchen mit kinetischer Energie $T=\SI{1000}{MeV}$ gilt
                    \[
                        \gamma=\frac{T+mc^2}{mc^2},
                        \qquad
                        \beta=\sqrt{1-\frac{1}{\gamma^2}}.
                    \]
                    Damit ergeben sich die folgenden Werte:
                    \[
                    \begin{array}{c|c|c|c|c|c}
                        \text{Teilchen} & z & \gamma & \beta
                        & -\dd E/\dd x\,[\si{MeV/cm}] & \Delta E\,[\si{MeV}]\\
                        \hline
                        p & 1 & 2.0658 & 0.8750 & 4.35 & 8.70\\
                        D & 1 & 1.5332 & 0.7580 & 5.39 & 10.78\\
                        \alpha & 2 & 1.2683 & 0.6151 & 30.61 & 61.21
                    \end{array}
                    \]
                    Dabei wurde jeweils
                    \[
                        \Delta E=\left(-\frac{\dd E}{\dd x}\right)d
                    \]
                    mit $d=\SI{2}{cm}$ verwendet.
                \end{calculationbox}
            \end{calculationbox}

            \begin{answerbox}
                Die Energieverluste im \SI{2}{cm} dicken Graphitabsorber sind näherungsweise
                \[
                    \boxed{\Delta E_p=\SI{8.70}{MeV},\qquad
                    \Delta E_D=\SI{10.78}{MeV},\qquad
                    \Delta E_{\alpha}=\SI{61.21}{MeV}.}
                \]
                Das $\alpha$-Teilchen verliert deutlich mehr Energie, hauptsächlich wegen des
                Faktors $z^2=4$ und wegen seiner kleineren Geschwindigkeit bei gleicher
                kinetischer Energie.
            \end{answerbox}

        \item
            Skizzieren Sie qualitativ den Energieverlust $-\dd E/\dd x$ über den Energiebereich
            von $\SI{1}{MeV}$ bis $\SI{1}{TeV}$. Nutzen Sie zur Normierung die Ergebnisse aus
            (a) und verwenden Sie logarithmische Achsen.

            \begin{calculationbox}
                Qualitativ besitzt die Bethe-Bloch-Kurve drei wichtige Bereiche:
                \begin{enumerate}[label=\arabic*.]
                    \item Bei kleinen Energien ist $\beta$ klein. Wegen des Faktors $1/\beta^2$
                    ist der Energieverlust sehr groß.
                    \item Bei mittleren Energien fällt die Kurve bis zu einem Minimum ab. Dieses
                    Gebiet nennt man den Bereich minimal ionisierender Teilchen.
                    \item Bei sehr hohen Energien wächst der Energieverlust durch den logarithmischen
                    Term langsam wieder an. Das ist der relativistische Anstieg.
                \end{enumerate}

                \begin{calculationbox}
                    Eine qualitative log-log-Skizze ist:
                    \begin{center}
                        \begin{tikzpicture}
                            \begin{loglogaxis}[
                                width=0.86\textwidth,
                                height=6.5cm,
                                xlabel={$T$ [MeV]},
                                ylabel={$-\dd E/\dd x$ [rel. Einheiten]},
                                xmin=1, xmax=1e6,
                                ymin=0.7, ymax=80,
                                grid=both,
                                legend pos=south west
                            ]
                            \addplot[thick,domain=1:1e6,samples=200]
                                {1 + 8/(x^(0.55)) + 0.20*ln(1+x/1000)};
                            \addlegendentry{qualitativ}
                            \addplot[only marks,mark=*] coordinates {(1000,4.35)};
                            \addlegendentry{$p$ bei \SI{1}{GeV}}
                            \addplot[only marks,mark=square*] coordinates {(1000,5.39)};
                            \addlegendentry{$D$ bei \SI{1}{GeV}}
                            \addplot[only marks,mark=triangle*] coordinates {(1000,30.61)};
                            \addlegendentry{$\alpha$ bei \SI{1}{GeV}}
                            \end{loglogaxis}
                        \end{tikzpicture}
                    \end{center}
                    Die eingezeichnete Kurve ist qualitativ. Die Punkte bei $\SI{1}{GeV}$ dienen
                    nur zur Normierung mit den Ergebnissen aus Teil (a).
                \end{calculationbox}
            \end{calculationbox}

            \begin{answerbox}
                Die Kurve fällt bei kleinen Energien stark ab, erreicht ein Minimum und steigt
                bei hohen Energien nur langsam logarithmisch wieder an. Die in Teil (a)
                berechneten Werte legen die Größenordnung für Proton, Deuteron und
                $\alpha$-Teilchen bei $\SI{1}{GeV}$ fest.
            \end{answerbox}

        \item
            Lässt sich die Bethe-Bloch-Formel auch für den Durchgang von Elektronen durch
            Materie verwenden? Begründen Sie Ihre Antwort.

            \begin{calculationbox}
                Die oben angegebene Bethe-Bloch-Formel ist in dieser einfachen Form für schwere
                geladene Teilchen gedacht, also zum Beispiel Myonen, Protonen, Deuteronen oder
                $\alpha$-Teilchen. Für Elektronen ist die Situation anders.

                \begin{calculationbox}
                    Der wichtigste Grund ist, daß das einlaufende Teilchen dann dieselbe Masse
                    hat wie die gebundenen Target-Elektronen. Deshalb sind die Annahmen der
                    Herleitung für schwere Projektile nicht mehr direkt erfüllt. Insbesondere
                    muß man die Ununterscheidbarkeit von Elektronen im Endzustand und die
                    veränderte Kinematik der Energieübertragung berücksichtigen.
                \end{calculationbox}

                \begin{calculationbox}
                    Zusätzlich werden bei Elektronen schon bei vergleichsweise niedrigeren
                    Energien Strahlungsverluste durch Bremsstrahlung wichtig. Für Elektronen
                    beschreibt man den Energieverlust daher mit modifizierten Formeln, die
                    Kollisionsverluste und Strahlungsverluste getrennt bzw. gemeinsam behandeln.
                \end{calculationbox}
            \end{calculationbox}

            \begin{answerbox}
                In der angegebenen einfachen Form sollte man die Bethe-Bloch-Formel nicht direkt
                für Elektronen verwenden. Elektronen benötigen eine modifizierte Behandlung,
                weil Projektil und Target-Elektronen gleiche Masse haben und weil Bremsstrahlung
                für sie eine deutlich wichtigere Rolle spielt.
            \end{answerbox}
    \end{enumerate}

\end{document}
